% Addition to root-owned NN-S016; no separate finding id.
\begin{proposition}[NN-S016: linear-degree integral coordinates in the empty case]
For the Boolean and verifier equations of NN-C027, suppose $N=0$.
There are integral polynomials $a_i(x)$ satisfying
\[
                      \sum_{i=1}^h a_i(x)f_i(x)=1,
 \qquad \deg a_i\leq v.
\]
Their coefficients have $O(v)$ bits. In the polynomial isomorphism
of the topology addendum, the forward coordinate polynomials have
degree at most $v+2$ and the inverse coordinate polynomials have
degree at most $v+6$, with coefficient bit lengths $O(v+\log(h+1))$.
These are degree and height bounds; a fully expanded coordinate
list can still have exponentially many monomials.
\end{proposition}

\begin{proof}
Every equation has degree at most two and coefficient $\ell^1$-norm
at most three. All $v$ Boolean equations $B_j=x_j^2-x_j$ are retained.
Every non-Boolean verifier equation takes a value in $\{0,1,-1\}$
on a bit assignment. Since there are no satisfying assignments,
for each $b\in\{0,1\}^v$ choose an index $i(b)$ with
$f_{i(b)}(b)\in\{1,-1\}$. Put
\[
 e_b(x)=\prod_{j:b_j=1}x_j\prod_{j:b_j=0}(1-x_j),
 \qquad
 a_i^{(0)}(x)=\sum_{b:i(b)=i}f_i(b)e_b(x).
\]
Then $\deg a_i^{(0)}\leq v$, and evaluation on the Boolean cube gives
$\sum_i a_i^{(0)}(b)f_i(b)=1$. In particular
\[
 P(x)=\sum_i a_i^{(0)}(x)f_i(x)-1
\]
has degree at most $v+2$ and has zero multilinear remainder modulo
the $B_j$. This last assertion holds over $\mathbb Z$: the
multilinear polynomial with prescribed values on the Boolean cube
is uniquely $\sum_b P(b)e_b$, so all its values being zero implies
that it is the zero polynomial.

Here is explicit integral division with controlled degree and
height. For an exponent vector $\alpha$, write
$\epsilon_j=\min(\alpha_j,1)$. Telescope
\[
 x^\alpha-x^\epsilon=
 \sum_j\left(\prod_{l<j}x_l^{\epsilon_l}\right)
       (x_j^{\alpha_j}-x_j^{\epsilon_j})
       \left(\prod_{l>j}x_l^{\alpha_l}\right).
\]
The middle factor is zero for $\alpha_j\leq1$, and otherwise it is
\[
        (x_j^2-x_j)\sum_{r=0}^{\alpha_j-2}x_j^r.
\]
Apply this formula to every monomial of $P$. Since its multilinear
remainder is zero, the result is
\[
                         P=\sum_j q_j B_j,
 \qquad \deg q_j\leq v.
\]
Subtract $q_j$ from the multiplier of the corresponding retained
Boolean equation. The resulting $a_i$ have the asserted degree
and satisfy the ordinary polynomial identity $\sum_i a_i f_i=1$.
This is an identity in the unrestricted integral polynomial ring,
not just in its Boolean quotient.

For a polynomial $Q$, write $\|Q\|_1$ for the sum of the absolute
values of its coefficients. Since $\|e_b\|_1=2^{|\{j:b_j=0\}|}$,
\[
 \sum_i\|a_i^{(0)}\|_1\leq\sum_b\|e_b\|_1=3^v,
 \qquad \|P\|_1\leq3^{v+1}+1.
\]
The division formula introduces at most $|\alpha|\leq v+2$
quotient monomials per input monomial, counted with multiplicity.
Hence
\[
 \sum_j\|q_j\|_1\leq(v+2)(3^{v+1}+1),\qquad
 \|a_i\|_1\leq B_v:=3^v+(v+2)(3^{v+1}+1).
\]
Thus every integer coefficient has $O(v)$ bits.

For the forward isomorphism $z_i=y_i+a_i(2t-t^2)$, the degree is
at most $v+2$ and its coefficient norm is at most $1+3B_v$.
For the inverse, $t=\sum_i f_i(x)z_i$ has degree at most three and
coefficient norm at most $3h$. Therefore
\[
 y_i=z_i+a_i\left(\left(\sum_j f_jz_j\right)^2
                           -2\sum_j f_jz_j\right)
\]
has degree at most $v+6$ and coefficient norm at most
$1+B_v(9h^2+6h)$. These give the claimed bit bounds.
A degree-$v$ polynomial in $v$ variables can have
$\binom{2v}{v}\leq4^v$ monomials; neither the degree nor the
coefficient-height estimate bounds their count polynomially.
\end{proof}

\paragraph{Proof-system and geometric interpretation.}
The multipliers form an integral Nullstellensatz refutation of the
quadratic verifier equations with their Boolean axioms. They also
supply low-degree polynomial coordinates for the same smooth cubic
model's arithmetic affine-space isomorphism. This refines the
Boolean-quotient identity of NN-C009 by explicitly constructing the
ordinary ideal multipliers in this special empty Boolean-scheme
case. No polynomial-size expanded certificate, short arithmetic
circuit, or efficient isomorphism-finding algorithm follows from
this degree calculation alone. NN-C029 separately analyzes what E's
symbolic partition routine supplies for circuit descriptions.
